% APF v38 -- Charged-Lepton Generation Residual Operator Theorem
\section{Charged-Lepton Generation Residual Operator}
\label{sec:charged-lepton-residual-operator-v38}

\paragraph{Purpose.}
The trace-to-scheme export theorem leaves the charged-lepton branch at a residual obstruction: after the scalar normalization candidate
\[
        \nu_\ell={1\over\sqrt 6}
\]
the branch still requires a generation-resolving residual operator.  This section closes the \emph{operator-form} part of that obstruction.  It proves that the missing object cannot be another scalar normalization, identifies the unique APF-minimal operator family on the three-generation charged-lepton interface, and records the remaining coefficient-closure gate.  The result is a route-object advance, not a physical-final export of electron, muon, and tau masses.

\subsection{Charged-lepton route data}

\begin{definition}[Charged-lepton trace vector]
The charged-lepton APF trace anchor is a three-component vector
\[
        T_\ell=(T_e,T_\mu,T_\tau)\in V_\ell,
        \qquad V_\ell\cong\mathbb R^3,
\]
produced inside the trace-sector generation algebra.  In the present package the vector is treated as an internal trace anchor: its component values are not imported from external electron, muon, or tau masses.
\end{definition}

\begin{definition}[Scalar normalization candidate]
The already-identified scalar normalization is
\[
        \nu_\ell={1\over\sqrt{3!}}={1\over\sqrt 6}.
\]
It accounts for the six admissible orderings of a three-generation branch before a generation orientation is selected.  Being scalar, it acts only on the uniform trace mode and cannot distinguish the three charged-lepton generations.
\end{definition}

\begin{definition}[Generation basis and trace-free modes]
Let
\[
        \mathbf 1=(1,1,1),\qquad q=(0,1,2),\qquad \bar q=1.
\]
The scalar normalization consumes the uniform mode spanned by \(\mathbf 1\).  The residual operator must therefore act on the trace-free generation plane
\[
        V_\ell^0=\{x\in V_\ell: x_e+x_\mu+x_\tau=0\}.
\]
APF supplies the canonical first-difference and second-difference generators
\[
        H_1:=\operatorname{diag}(-1,0,1),
        \qquad
        H_2:=\operatorname{diag}(1,-2,1),
\]
which span all diagonal trace-free generation operators compatible with the ordered three-rung ladder.
\end{definition}

\begin{definition}[Forbidden charged-lepton residual inputs]
A charged-lepton residual operator is APF-native only if neither its form nor its coefficients are chosen using
\[
        \mathcal F_\ell=
        \{m_e^{\rm pole},m_\mu^{\rm pole},m_\tau^{\rm pole},
        m_e(\mu),m_\mu(\mu),m_\tau(\mu),
        \hbox{PDG/world-average charged-lepton masses},
        \hbox{least-squares fits to the charged-lepton target vector}\}.
\]
Those quantities may be used later for external comparison, but not to construct the APF residual route.
\end{definition}

\subsection{Scalar insufficiency and the minimal operator}

\begin{lemma}[Scalar normalization cannot close a three-generation residual]
\label{lem:scalar-insufficiency-v38}
No scalar multiple of \(\nu_\ell T_\ell\) can be an APF-complete charged-lepton export route unless the residual vector is generation-blind.
\end{lemma}

\begin{proof}
A scalar normalization acts as \(a I\) on \(V_\ell\).  It changes only the overall scale of the vector and commutes with every permutation of the three generations.  Therefore it cannot assign different residual load to the first, second, and third generation.  But a charged-lepton export target is generation-resolved: it distinguishes \(e\), \(\mu\), and \(\tau\).  Once the uniform ordering factor \(1/\sqrt6\) is placed on-ledger, any remaining generation-dependent mismatch lies in \(V_\ell^0\).  A scalar has zero action on the distinction between directions in \(V_\ell^0\).  Hence a non-scalar generation operator is necessary unless the residual is generation-blind, which would erase the charged-lepton hierarchy rather than export it. \qedhere
\end{proof}

\begin{theorem}[APF-minimal charged-lepton residual operator]
\label{thm:charged-lepton-residual-operator-form-v38}
Under the APF requirements of generation typing, closed scalar normalization, ordered three-rung FN ladder compatibility, diagonal charged-lepton codomain admission, and no target reuse, the minimal charged-lepton residual operator has the form
\[
        \mathcal R_\ell(\alpha,\beta)
        =\exp\!\left(\alpha H_1+\beta H_2\right),
\]
where
\[
        H_1=\operatorname{diag}(-1,0,1),
        \qquad
        H_2=\operatorname{diag}(1,-2,1).
\]
Thus the admissible route object is
\[
        \Phi_\ell:
        T_\ell\longmapsto
        O_\ell^S
        =Z_\ell^S\,\nu_\ell\,
        \mathcal R_\ell(\alpha,\beta)T_\ell,
\]
where \(S\) is a declared charged-lepton codomain and \(Z_\ell^S\) is the scheme-specific QED/running factor if the codomain is not pole-normalized.
\end{theorem}

\begin{proof}
The residual must be an endomorphism of the charged-lepton generation space, because the source is the trace vector \(T_\ell\in V_\ell\) and the target distinguishes the same three charged-lepton generations.  The scalar ordering normalization \(\nu_\ell=1/\sqrt6\) has already been assigned to the uniform mode, so any genuine residual must be trace-free on generation space.  In the charged-lepton mass basis no off-diagonal mixing operator is admissible unless an additional PMNS/charged-current transport map is supplied; otherwise the route would silently import neutrino-sector structure.  Therefore the minimal residual is diagonal and trace-free.

A diagonal trace-free operator on three ordered generations is two-dimensional.  The APF generation ladder supplies the ordered charges \((0,1,2)\), hence the natural polynomial basis after removing the mean is the first-difference mode and the second-difference mode:
\[
        H_1=\operatorname{diag}(-1,0,1),
        \qquad
        H_2=\operatorname{diag}(1,-2,1).
\]
Every diagonal trace-free generation operator is a unique linear combination \(\alpha H_1+\beta H_2\).  Exponentiation is forced by positive mass transport and composition: sequential residual transports add in the generator and multiply on the positive mass vector.  Thus the APF-minimal positive residual operator is exactly
\[
        \mathcal R_\ell(\alpha,\beta)=\exp(\alpha H_1+\beta H_2).
\]
The coefficients \(\alpha\) and \(\beta\) are not fit parameters in an export claim.  They must be derived from APF generation-interface invariants or else the charged-lepton branch remains at a named coefficient-closure obstruction. \qedhere
\end{proof}

\begin{corollary}[Charged-lepton branch status after v38]
The charged-lepton obstruction is sharpened from ``missing generation residual operator'' to
\[
        \boxed{\text{operator form closed; coefficient closure still open}.}
\]
The branch may be assigned
\[
        \ell:[P_{\rm route}^{\rm operator\ form}],
\]
but not yet
\[
        \ell:[P_{\rm export\ candidate}],
        \qquad
        \ell:[P_{\rm physical\ final}].
\]
\end{corollary}

\begin{proof}
G1 is partially satisfied once a pole or running-QED codomain \(S\) is declared.  G2 is advanced because the transport map must be \(\Phi_\ell=Z_\ell^S\nu_\ell\mathcal R_\ell T_\ell\) with \(\mathcal R_\ell\) as above.  G3 records the internal ledger \(3!,\nu_\ell,q,H_1,H_2\) and any later QED factors.  G4 has the formal covariance propagation
\[
        \Sigma_{O_\ell}
        =D\Phi_\ell\,\Sigma_{(T_\ell,\alpha,\beta,Z)}D\Phi_\ell^{*}
        +\Sigma_{\rm QED}+\Sigma_{\rm scheme}.
\]
G5 passes only for the operator form, because the form uses no target masses.  G6 is sharpened: remaining residuals are coefficient closure, QED running, pole/running scheme choice, covariance, or APF route residual.  Since \(\alpha\) and \(\beta\) are not yet evaluated from APF-native data in this package, export-candidate status is not licensed. \qedhere
\end{proof}

\subsection{Coefficient-closure gate}

\begin{definition}[Charged-lepton coefficient-closure gate]
The remaining APF-native theorem must supply a source-independent map
\[
        \mathfrak C_\ell:
        \Gamma_{\rm gen}\longmapsto(\alpha_\ell,\beta_\ell)
\]
from generation-interface data alone.  Acceptable inputs include APF generation charges, finite ordering volume, cooperative-kernel invariants, and charged-lepton trace anchors.  Forbidden inputs are the external charged-lepton target masses and any least-squares residual fit to them.
\end{definition}

\begin{center}
\renewcommand{\arraystretch}{1.2}
\begin{tabular}{lll}
\hline
\textbf{Gate} & \textbf{v38 charged-lepton status} & \textbf{Reason}\\
\hline
G1 & partial & pole/running QED codomain must be declared for comparison\\
G2 & advanced & residual operator form \(\exp(\alpha H_1+\beta H_2)\) derived\\
G3 & partial & internal ledger closed; QED constants deferred until codomain choice\\
G4 & formal & covariance formula supplied, numeric covariance deferred\\
G5 & pass for form & no charged-lepton target masses used to derive the form\\
G6 & pass as obstruction & remaining channels named, coefficients still open\\
\hline
\end{tabular}
\end{center}

\paragraph{Bottom line.}
The charged-lepton branch is no longer blocked by an undefined type of missing object.  The missing route object has a forced APF-minimal form:
\[
        T_\ell\mapsto
        Z_\ell^S{1\over\sqrt6}
        \exp(\alpha H_1+\beta H_2)T_\ell.
\]
What remains is the harder coefficient theorem: derive \(\alpha\) and \(\beta\) from APF generation-interface structure without consuming the electron, muon, or tau target masses.
